\section{Related Work}
\label{sec:related_work}

\subsection{Low-Rank Tensor Recovery}

CP and Tucker decompositions are widely used for compression, denoising,
source separation, completion, and data fusion
\cite{Kolda2009Tensor,Cichocki2015SignalProcessing,Sidiropoulos2017TSP}.
Tucker decomposition represents a tensor by a core and mode-wise factors
\cite{DeLathauwer2000MLSVD}, so its multilinear ranks directly specify the
capacity assigned to each mode. Tensor-train and tensor-ring decompositions
control complexity through different network structures
\cite{oseledets2011,zhao2016tr}. These formats provide useful alternative
low-rank priors, while Tucker is the focus here because its mode-wise ranks let
us compare the rank of the latent tensor with the rank observed after a
nonlinear response.

Tensor completion methods recover missing entries by constraining the tensor
with convex surrogates or factorized low-rank models. Representative approaches
include nuclear-norm-based formulations, Tucker or CP factorizations, and
nonconvex algorithms for incomplete tensors
\cite{Liu2013TensorCompletion,LiuMoitra2020TensorCompletion,yang2016iterative}.
These methods generally place the low-rank constraint directly on the measured
tensor. In our setting, the Tucker constraint is placed on a latent
pre-response tensor, while the measured tensor can have a larger apparent
rank. Increasing the Tucker rank remains an important baseline, but it adds
parameters throughout the core and factors. NTD-PL instead tests whether the
additional measured structure can be explained by a shared scalar response at
the prescribed latent Tucker rank.

\subsection{Polynomial and Nonlinear Tensor Models}

Polynomial predictors can be written with high-order coefficient tensors.
Tensor factorizations have been used to compress these coefficients and avoid
forming all polynomial features explicitly
\cite{yang2015tensorMachines,novikov2016exponentialMachines}. These methods
produce scalar predictions from vector inputs. The Tucker polynomial considered
here has a different role: it maps a Tucker core to a tensor output, while the
mode maps preserve the structure of each output dimension. NTD-PL is the
shared-factor form of this model.

Generalized tensor factorizations combine low-rank structure with a prescribed
observation distribution or link. Generalized CP models support a range of
loss functions and data types, while probabilistic Tucker models use
application-specific likelihoods for observations such as counts
\cite{hong2020generalized,schein2016bayesian}. Nonparametric Bayesian tensor
models provide another form of flexibility through latent-function priors
\cite{xu2011infinite}. In these approaches, the observation model is typically
chosen as part of the statistical formulation.

Kernel, functional, and neural tensor models learn broader nonlinear feature
maps or decoders
\cite{Luo2024LRTFR,liu2017deepcp,saragadam2024deeptensor,ahn2024neural}.
They are designed to capture nonlinear interactions that need not reduce to a
single entrywise response. NTD-PL addresses a narrower case: one scalar
function is shared across entries, and the input to that function is a Tucker
tensor with explicitly prescribed ranks. This restriction retains a direct
connection between the latent rank budget and the complexity added in the
measured domain.

The closest matrix analogue is low-rank estimation with a generalized loss or
an unknown entrywise transfer function. Generalized low-rank models cover many
data-dependent losses, and monotonic single-index matrix completion jointly
estimates a low-rank matrix and an unknown monotone link
\cite{Udell2016GLRM,Ganti2015SingleIndex}. The present work follows the same
view for Tucker tensors: low-rank structure is modeled before an entrywise
response. The Tucker polynomial supplies the higher-order tensor map, while
sharing its factors yields a small response model that can be fitted and
diagnosed efficiently.

\subsection{Nonlinear Responses in Imaging}

Imaging pipelines provide established examples of nonlinear measurement
responses. Camera response functions map scene irradiance to recorded pixel
values, and high-dynamic-range imaging often estimates or calibrates this map
explicitly \cite{Debevec1997HDR,Grossberg2004CameraResponse}. These models
motivate learning a low-dimensional response instead of absorbing all
radiometric variation into the image representation.

Hyperspectral imaging contains other forms of nonlinearity, especially in
spectral mixing. Reviews of nonlinear unmixing and post-nonlinear mixing models
describe settings in which a structured linear mixture is followed by a
nonlinear transformation
\cite{BioucasDias2012HyperspectralUnmixing,Altmann2012PostNonlinear,
Heylen2014NonlinearUnmixing}. These physical mechanisms are not all equivalent
to one global entrywise response. Their relevance here is the broader modeling
precedent: a compact structured signal can be considered before a nonlinear
observation stage.

Accordingly, NTD-PL is not presented as a complete physical model of a camera
or hyperspectral sensor, and the real-data experiments do not claim to identify
a unique calibration curve. Instead, we test whether a jointly learned shared
response improves reconstruction and completion of the measured tensor at a
fixed latent Tucker rank. Hyperspectral cubes are useful for this test because
of their strong spatial and spectral correlations. The response diagnostic
further measures, before joint fitting, whether a scalar response can reduce
the residual of a Tucker model on a given tensor.
